\documentclass[minion]{homework}

% Some common macros
\newcommand{\Reals}{\ensuremath{\mathbb R}}% Gives you a shortcut for writing the blackboard R for the real numbers - \RR
\newcommand{\Rats}{\ensuremath{\mathbb Q}} % Gives you a shortcut for writing the blackboard N for the natural numbers - \NN
\newcommand{\Nats}{\ensuremath{\mathbb N}} % Gives you a shortcut for writing the blackboard N for the natural numbers - \NN
\newcommand{\Ints}{\ensuremath{\mathbb Z}} % Gives you a shortcut for writing the blackboard Z for the integer numbers - \ZZ
\doclabel{Christopher Munoz \qquad Math F404: Homework 3}
\docdate{Due: September 15, 2026}

\begin{document}
\begin{problems}

\problem  Suppose $(x_n)$ is a real sequence and $x_n\to \ell > 0$.
Show that there exists $N$ so that if $n\ge N$ then $x_n>\ell/2$.
Then show $1/x_n \to 1/\ell$. Note that the sequence $1/x_n$ is 
perhaps only defined for $n\ge N$ to avoid division by zero.
\newline
\begin{proof}[Part 1] $(x_n \to \ell > 0 \implies x_n > \ell/2) \because$ \newline
  Given $x_n \to \ell$, by definition there exists an $n \geq N$ such that $|x_n - \ell| < \varepsilon$ for any $\varepsilon > 0$. Note that $\ell/2 > 0$ and choose $\varepsilon = \ell/2$. Then 
  $$|x_n - \ell| < \frac{\ell}{2}$$
  holds, in particular we get $x_n > \ell - \ell/2$. Henceforth $x_n > \ell/2$.
\end{proof}$ $
\begin{proof}[Part 2] $(x_n > \ell/2 \implies 1 / x_n \to 1/\ell)\because$
  Let $\varepsilon > 0$. Since $x_n \to l$, by definition there exists $N' \geq N$ such that for all $n \geq N'$:
  $$|x_n - \ell| < \frac{\varepsilon\ell^2}{2}$$
  For $n \geq N'$(since $n \geq N, x_n > \ell / 2$ is given by Part 1):
      \begin{align*}
      \left|\frac{1}{x_n} - \frac{1}{\ell}\right| &= \frac{|\ell - x_n|}{x_n \ell} 
      < \frac{|\ell-x_n|}{(\ell/2)\cdot \ell} 
      = \frac{2|x_n-\ell|}{\ell^2} 
      < \frac{2}{\ell^2}\cdot\frac{\varepsilon\ell^2}{2}
      = \varepsilon.
    \end{align*}
    Thus $1/x_n \to 1/\ell$.
\end{proof}

\problem Suppose $\Omega\subset\Reals$, that $a$ is a limit point of $\Omega$, 
and that $f:\Omega\to\Reals$.  Show that $\lim_{x\to a}f(x) =\ell$ if
and only if whenever $(x_n)$ is a sequence in $\Omega\setminus\{a\}$ converging to $a$,
$f(x_n) \to \ell$.  You are encouraged to modify the similar proof seen in 
class on the sequential characterization of continuity.
\newline
\begin{proof} \textbf{Forward Direction} \newline
  Suppose $\lim_{x \to a} f(x) = \ell$. Let $(x_n)$ be a sequence in $\Omega \setminus \{a\}$ with $x_n \to a$. Let $\varepsilon > 0$. By the $\varepsilon-\delta$ definition, there exists $\delta > 0$ such that for all $x \in \Omega$,
  $$0 < |x - a| < \delta \implies |f(x) - \ell| < \epsilon.$$
  Since $x_n \to a$, there exists $N$ such that $|x_n - a| < \delta$ for all $n \geq N$. Since $x_n \neq a$ for every $n$ (giving $0 < |x_n - a|)$, for $n \geq N$, we have $0 < |x_n - a| < \delta$, so
  $$|f(x_n)-\ell| < \varepsilon$$
  Thus $f(x_n) \to \ell$.
\end{proof}$ $
\begin{proof} \textbf{Backwards Direction, by contraposition}
  Suppose $\lim_{x \to a} f(x) = \ell$ is false. Negating the $\varepsilon-\delta$ definition means there exists a $\varepsilon > 0$ such that
  for all $\delta > 0$, there exists an $x \in \Omega$ with
  $$0 < |x - a| < \delta \text{ and } |f(x) - \ell| \geq \varepsilon$$
  Fix this $\varepsilon$. For each $n \in \Nats$, let $\delta = 1/n$ to obtain $x_n \in \Omega$ with
  $$ 0 < |x_n - a| < \frac{1}{n} \text{  and  } |f(x_n) - \ell| \geq \varepsilon.$$
  Since $0 < |x_n - a|$, $x_n \in \Omega \setminus \{a\}$. Also since $|x_n - a| < 1/n \to 0$, $x_n \to a$. But $|f(x_n) - \ell| \geq \varepsilon$ for every $n$, so $f(x_n)$ does not converge to $\ell$.
  Thus $(x_n)$ is a sequence in $\Omega \setminus \{a\}$ with $x_n \to a$ but $f(x_n) \not\to \ell$, so the sequential condition fails. This proves the contrapositive, so the backward direction holds.
\end{proof}
\problem Sutherland 4.13 \newline
(a) Prove that for any $y,z\in\Reals$,
\[
\max\{y,z\} = \frac12(y+z+|y-z|), \qquad \min\{y,z\} = \frac12(y+z-|y-z|).
\]
(b) Given that $f,g:\Reals\to\Reals$ are continuous at $a$, prove that $h$ and $k$ are
continuous at $a\in\Reals$, where for any $x\in\Reals$
\[
h(x)=\max\{f(x),g(x)\}, \qquad k(x)=\min\{f(x),g(x)\}.
\]
\newline
\begin{proof}[(a)]
  Let $y,z\in\Reals$. We consider two cases.

  \textbf{Case 1: $y\ge z$.} Then $y-z\ge0$, so $|y-z|=y-z$. Thus
  $$\frac12(y+z+|y-z|) = \frac12(y+z+y-z) = y = \max\{y,z\},$$
  $$\frac12(y+z-|y-z|) = \frac12(y+z-y+z) = z = \min\{y,z\}.$$

  \textbf{Case 2: $y<z$.} Then $y-z<0$, so $|y-z|=z-y$. Thus
  $$\frac12(y+z+|y-z|) = \frac12(y+z+z-y) = z = \max\{y,z\},$$
  $$\frac12(y+z-|y-z|) = \frac12(y+z-z+y) = y = \min\{y,z\}.$$

  In either case the two identities hold.
\end{proof}$ $
\begin{proof}[(b)]
  By part (a), for every $x\in\Reals$,
  $$h(x) = \frac12\big(f(x)+g(x)+|f(x)-g(x)|\big), \qquad k(x) = \frac12\big(f(x)+g(x)-|f(x)-g(x)|\big).$$
  Let $(x_n)$ be any sequence with $x_n\to a$. Since $f,g$ are continuous at $a$, by the sequential characterization of continuity, $f(x_n)\to f(a)$ and $g(x_n)\to g(a)$. By the limit laws for sequences,
  $$f(x_n)+g(x_n) \to f(a)+g(a), \qquad f(x_n)-g(x_n)\to f(a)-g(a).$$
  Since $\big||u|-|v|\big|\le|u-v|$ for all $u,v\in\Reals$ (the absolute value function is continuous), applying this with $u=f(x_n)-g(x_n)$ and $v=f(a)-g(a)$ gives
  $$\big||f(x_n)-g(x_n)| - |f(a)-g(a)|\big| \le |(f(x_n)-g(x_n))-(f(a)-g(a))| \to 0,$$
  so $|f(x_n)-g(x_n)|\to|f(a)-g(a)|$. Combining these limits,
  $$h(x_n) = \frac12\big(f(x_n)+g(x_n)+|f(x_n)-g(x_n)|\big) \to \frac12\big(f(a)+g(a)+|f(a)-g(a)|\big) = h(a),$$
  $$k(x_n) = \frac12\big(f(x_n)+g(x_n)-|f(x_n)-g(x_n)|\big) \to \frac12\big(f(a)+g(a)-|f(a)-g(a)|\big) = k(a).$$
  Since $(x_n)$ was an arbitrary sequence converging to $a$, $h$ and $k$ are continuous at $a$.
\end{proof}

\problem Sutherland 4.1 \textbf{(AI Free)} \newline
Let $f,g:\Reals\setminus\{0\}\to\Reals$ be given by
\[
f(x) = x\sin\frac1x, \qquad g(x) = \sin\frac1x.
\]
Prove that $\lim_{x\to0} f(x) = 0$, while $\lim_{x\to0} g(x)$ does not exist.
\newline
\begin{proof}[$f$]
  Let $\varepsilon>0$ and choose $\delta=\varepsilon$. Then for $0<|x|<\delta$, since $|\sin(1/x)|\le1$,
  $$|f(x) - 0| = |x||\sin(1/x)| \le |x| < \delta = \varepsilon.$$
  Hence $\lim_{x\to0} f(x) = 0$.
\end{proof}$ $
\begin{proof}[$g$]
  Suppose for contradiction that $\lim_{x\to0} g(x) = \ell$ for some $\ell\in\Reals$. By the sequential characterization of limits (Problem 2, with $\Omega=\Reals\setminus\{0\}$ and $a=0$), every sequence $(x_n)$ in $\Omega\setminus\{0\}$ with $x_n\to0$ must satisfy $g(x_n)\to\ell$.

  Let $x_n = \dfrac1{2n\pi+\pi/2}$ and $y_n = \dfrac1{2n\pi-\pi/2}$ for $n\in\Nats$. Both sequences lie in $\Reals\setminus\{0\}$ and $x_n,y_n\to0$. But
  $$g(x_n) = \sin\!\left(2n\pi+\frac\pi2\right) = 1 \quad\text{for all }n, \qquad g(y_n) = \sin\!\left(2n\pi-\frac\pi2\right) = -1 \quad\text{for all }n.$$
  So $g(x_n)\to1$ and $g(y_n)\to-1$. By the sequential characterization, both limits must equal $\ell$, forcing $1=\ell=-1$, a contradiction. Hence $\lim_{x\to0} g(x)$ does not exist.
\end{proof}

\problem Sutherland 4.15
Let $f:\Reals\to\Reals$ be given by
\[
f(x) = \begin{cases} 0, & x\in\Rats \\ 1, & x\notin\Rats \end{cases}
\]
Show that $f$ is not continuous at any point in $\Reals$.
\newline
\begin{proof}
  Let $a\in\Reals$ be arbitrary. By density of $\Rats$ in $\Reals$, there exists a sequence $(p_n)$ of rationals with $p_n\to a$; by density of $\Reals\setminus\Rats$ in $\Reals$, there exists a sequence $(q_n)$ of irrationals with $q_n\to a$. By definition of $f$,
  $$f(p_n) = 0 \text{ for all } n, \qquad f(q_n) = 1 \text{ for all } n,$$
  so $f(p_n)\to0$ and $f(q_n)\to1$.

  If $a\in\Rats$, then $f(a)=0$, but $q_n\to a$ while $f(q_n)\to1\neq0=f(a)$, so by the sequential characterization of continuity, $f$ is not continuous at $a$.

  If $a\notin\Rats$, then $f(a)=1$, but $p_n\to a$ while $f(p_n)\to0\neq1=f(a)$, so again $f$ is not continuous at $a$.

  Since $a$ was arbitrary, $f$ is not continuous at any point of $\Reals$.
\end{proof}

\problem Use Definition 4.21 and no more advanced tools to directly show 
\[
\lim_{x\to 2} \frac{1}{x^2+1} = \frac{1}5.
\]
That is, start your proof with "Let $\epsilon>0$", exhibit a candidate
$\delta$, and show that your candidate works.
\newline
\begin{proof}
  Let $\varepsilon>0$. For $x\in\Reals$,
  $$\left|\frac1{x^2+1} - \frac15\right| = \left|\frac{5-(x^2+1)}{5(x^2+1)}\right| = \frac{|4-x^2|}{5(x^2+1)} = \frac{|x-2||x+2|}{5(x^2+1)}.$$
  Suppose $|x-2|<1$, so $1<x<3$. Then $x^2+1>2$ and $0<x+2<5$, so $|x+2|<5$. Hence
  $$\left|\frac1{x^2+1}-\frac15\right| = \frac{|x-2||x+2|}{5(x^2+1)} < \frac{|x-2|\cdot5}{5\cdot2} = \frac{|x-2|}2.$$
  Let $\delta = \min\{1, 2\varepsilon\}$. Then for $0<|x-2|<\delta$, we have $|x-2|<1$, so the bound above applies, and $|x-2|<2\varepsilon$, so
  $$\left|\frac1{x^2+1}-\frac15\right| < \frac{|x-2|}2 < \frac{2\varepsilon}2 = \varepsilon.$$
  Hence $\lim_{x\to2}\dfrac1{x^2+1} = \dfrac15$.
\end{proof}

\problem For $x,y\in \Reals$ define $d(x,y) = |y-x|$.
Show that a sequence of real numbers $(x_n)$ converges to
a limit $\ell$ if and only if $d(x_n, \ell)\to 0$.
\newline
\begin{proof}
  Note that for every $n$, $d(x_n,\ell) = |\ell - x_n| = |x_n - \ell|$.

  ($\implies$) Suppose $x_n\to\ell$: for every $\varepsilon>0$ there exists $N$ such that $n\ge N$ implies $|x_n-\ell|<\varepsilon$. Since $d(x_n,\ell) = |x_n-\ell|$, this says $n\ge N$ implies $|d(x_n,\ell) - 0| < \varepsilon$, i.e., $d(x_n,\ell)\to0$.

  ($\impliedby$) Suppose $d(x_n,\ell)\to0$: for every $\varepsilon>0$ there exists $N$ such that $n\ge N$ implies $|d(x_n,\ell)-0|<\varepsilon$. Since $d(x_n,\ell) = |x_n-\ell|$, this says $n\ge N$ implies $|x_n-\ell|<\varepsilon$, i.e., $x_n\to\ell$.

  Hence $x_n\to\ell$ if and only if $d(x_n,\ell)\to0$.
\end{proof}

\problem Sutherland 5.1
Given points $x,y,z$ in a metric space $(X,d)$, prove that
\[
|d(x,z) - d(y,z)| \le d(x,y).
\]
\begin{proof}
  By the triangle inequality, $d(x,z) \le d(x,y) + d(y,z)$, so
  $$d(x,z) - d(y,z) \le d(x,y).$$
  Also by the triangle inequality and symmetry of $d$, $d(y,z) \le d(y,x)+d(x,z) = d(x,y)+d(x,z)$, so
  $$d(y,z) - d(x,z) \le d(x,y), \quad\text{i.e.,}\quad -(d(x,z)-d(y,z)) \le d(x,y).$$
  Combining the two inequalities, $-d(x,y) \le d(x,z)-d(y,z) \le d(x,y)$, which is equivalent to
  $$|d(x,z)-d(y,z)| \le d(x,y).$$
\end{proof}




\end{document}
